This section develops a comprehensive framework for analyzing stability and
performance of interconnected systems in a Hilbert space, utilizing IQCs as the
primary analytical tool.
We begin by establishing the mathematical foundation for system interconnections
and formalizing the key concepts of stability and performance.
The IQC methodology is then systematically developed to characterize system
uncertainties and derive rigorous conditions that guarantee robust stability and
performance.
A central contribution of this work lies in bridging IQC theory with
$\mu$-analysis for infinite dimensional systems, creating a framework for
quantifying their robust stability and performance when those systems are
subjected to structured uncertainties.

\subsection{Description of uncertain systems in a Hilbert-space}
We start by denoting the uncertainty operators as $\Delta$ that belongs to a
structured set $\mathbf{\Delta}$ which satisfies the following fundamental
property:
\begin{assumption}
$\Delta\in\mbf{\Delta}$ implies that $\tau\Delta\in\mbf{\Delta}$ for all
$\tau \in [0,1]$.
\end{assumption}
That is, the uncertainty set is closed under scaling by factors in $[0,1]$,
which is the standard homotopy assumption used in IQC theory.
We consider operators that map external inputs to external outputs through a
feedback interconnection between a linear (infinite dimensional) operator $P$
and $\Delta$.
This interconnection structure is formally defined as follows
\begin{definition}\label{def:uncertain-interconnection}
Given operators
$P_\Delta :
\mbf L_{e,[a,b]}^{\mbf{n}_{w_\Delta}}
\to
\mbf L_{e,[a,b]}^{\mbf{n}_{z_\Delta}}$,
$P_{\Delta p} :
\mbf L_{e,[a,b]}^{\mbf{n}_{w}}
\to
\mbf L_{e,[a,b]}^{\mbf{n}_{z_\Delta}}$,
$P_{p\Delta} :
\mbf L_{e,[a,b]}^{\mbf{n}_{w_\Delta}}
\to
\mbf L_{e,[a,b]}^{\mbf{n}_{z}}$,
$P_{p} :
\mbf L_{e,[a,b]}^{\mbf{n}_{w}}
\to
\mbf L_{e,[a,b]}^{\mbf{n}_{z}}$,
and $\Delta :
\mbf L_{e,[a,b]}^{\mbf{n}_{z_\Delta}}
\to
\mbf L_{e,[a,b]}^{\mbf{n}_{w_\Delta}}$, define
\begin{equation}\label{eq:uncertain-interconnection-map}
\bmat{z_\Delta\\ z_p}
=
\bmat{
P_\Delta & P_{\Delta p} \\
P_{p\Delta} & P_{p}
}\bmat{w_\Delta\\ w_p},
\quad
w_\Delta=\Delta z_\Delta.
\end{equation}
for input $w_p \in \mbf L_{e,[a,b]}^{\mbf{n}_{w}}$ and output
$z_p\in \mbf L_{e,[a,b]}^{\mbf{n}_{z}}$, there exist internal signals
$z_\Delta \in \mbf L_{e,[a,b]}^{\mbf{n}_{z_\Delta}}$ and
$w_\Delta \in \mbf L_{e,[a,b]}^{\mbf{n}_{w_\Delta}}$
satisfying the {interconnected system} $z_p= (P\star\Delta)w_p$.
\end{definition}
The interconnection structure described above captures the feedback interaction
between the nominal system $P$ and the uncertainty $\Delta$, where $w_p$ and $z_p$
represent external inputs and performance outputs, respectively, while
$w_\Delta$ and $z_\Delta$ form the internal feedback loop through the
uncertainty.
Before analyzing stability and performance, we must first ensure the
well-posedness of the interconnection, such that solutions exist, are unique,
and depend causally on inputs.
\begin{definition}
We say that the uncertain interconnection $P\star\Delta$ is
\textbf{well--posed} if:
\begin{itemize}
\item \textbf{Existence and Uniqueness}: For any
$w\in \mbf L_{e, [a, b]}$, there exist unique
$z\in\mbf L_{e, [a, b]}$,
$z_\Delta \in \mbf L_{e, [a, b]}$ and
$w_\Delta \in  \mbf L_{e, [a, b]}$ such that, $z$, $w$,
$z_\Delta, w_\Delta$ satisfy the interconnection.
\item \textbf{Causality}: If $z_\Delta, w_\Delta$ and
$\hat{z}_\Delta, \hat{w}_\Delta$ satisfy the interconnection,
then $p_T(z_\Delta - \hat{z}_\Delta) = 0$ and
$p_T(w_\Delta - \hat{w}_\Delta) = 0$, for all
$T \geq 0$.
\end{itemize}
\end{definition}
Having established the class of uncertain systems and their well-posedness, we
now develop a framework to guarantee robust stability and performance for the
class of systems defined in Definition~\ref{def:uncertain-interconnection}.

\subsection{Robust stability and performance}
Traditional control system design often treats stability and performance
shaping as separate objectives with distinct constraints.
However, Integral Quadratic Constraints (IQCs) provide a unified framework where
these trade-offs can be captured within a single constraint \cite{10156335}.
We now define what it means for the system to remain stable despite
uncertainties.
Robust stability requires that the system is well-posed for all admissible
uncertainties and that the input-output map remains bounded.
This involves removal of the performance channels in
\eqref{eq:uncertain-interconnection-map}, and
introducing auxiliary signals $z_0$ and $w_0$.
\begin{definition}
The interconnection
\begin{align*}
z_\Delta = P_\Delta w_\Delta + z_0, \quad
w_\Delta = \Delta z_\Delta + w_0,
\end{align*}
with $z_0, w_0 \in \mbf{L}_{e,[a,b]}$, is said to be
\textbf{robustly stable} if:
\begin{enumerate}
\item The uncertain system is well-posed for all $\Delta\in\mbf{\Delta}$.
\item The map from $\bmat{z_0\\w_0}$ to
$\bmat{z_\Delta\\w_\Delta}$ is bounded on $\mbf L_{[a,b]}$
for all $\Delta\in\mbf{\Delta}$.
\end{enumerate}
\end{definition}
Within the IQC framework, the stability of an interconnection $P\star\Delta$ is
certified by verifying that the uncertainty operator $\Delta$ satisfies a
quadratic constraint defined by an appropriate multiplier.
This leads to the following formal definition
\begin{definition}\label{def:uncertainty-iqc}
An uncertainty operator
$\Delta:
\mbf L_{e,[a,b]}^{\mbf{n}_{z_\Delta}}
\to
\mbf L_{e,[a,b]}^{\mbf{n}_{w_\Delta}}$
is said to satisfy the IQC defined by the operator
$\Psi_\Delta\in
\mbf L^{\mbf{n}_{z_\Delta} + \mbf{n}_{w_\Delta}}_{e, [a, b]}
\to
\mbf L^{\mbf{n}_{z_\Delta} + \mbf{n}_{w_\Delta}}_{e, [a, b]}$
multiplier $\mcl M_{\mcl V_\Delta}\in \mcl L(\mbf L_{[a, b]})$,
parameterized by
$\mcl{V}_\Delta = \mcl{V}_\Delta^* \in \mcl L(\RL, \RL)$
if the inequality holds for all
$z_\Delta \in \mbf L_{e, [a, b]}$ and all $T \geq 0$,
\begin{equation}\label{eq:uncertainty-iqc}
\ip{p_T \Psi_\Delta \bmat{z_\Delta \\ w_\Delta}}
{p_T \mcl{M}_{\mcl{V}_\Delta}\Psi_\Delta\bmat{z_\Delta \\ w_\Delta}}
_{\mbf L_{[a, b]}}
\geq 0.
\end{equation}
Moreover, let
$\mcl{V}_\Delta\{\mcl{Q}_\Delta, \mcl{S}_\Delta, \mcl{R}_\Delta\}$
be parameterized as,
\begin{equation*}
\mcl{V}_\Delta =
\bmat{
\mcl{Q}_\Delta & \mcl{S}_\Delta \\
\mcl{S}_\Delta^* & \mcl{R}_\Delta
}
\end{equation*}
\end{definition}
Beyond stability, we are interested in ensuring that the system meets certain
performance specifications in the presence of uncertainty.
Robust performance extends the notion of robust stability by additionally
requiring that a performance metric is satisfied.
\begin{definition}\label{def:robust-performance}
\textbf{Robust performance} is achieved if the interconnection $P \star \Delta$
is \textbf{robustly stable} for some performance multiplier
$\mcl {M}_{{\mcl {V}_p}}\in\mcl L(\mbf L_{[a, b]})$,
parameterized by
${\mcl{V}}_p = {\mcl{V}}^*_p\in \mcl L(\RL, \RL)$
with ${\mcl Q}_{p}\succeq 0$ such that, for all
$w_p \in\mbf L_{e, [a, b]}$, all $T \geq 0$, and some
sufficiently small $\epsilon >0$,
\begin{equation*}
\ip{p_T\bmat{z_p \\ w_p}}
{p_T\mcl{M}_{{\mcl{V}}_p}\bmat{z_p \\ w_p}}_{\mbf L_{[a, b]}}
\!\!\!\!\!\!\!\!\!\!\!\!\leq
-\epsilon \|p_Tw_p\|_{\mbf L_{[a, b]}}^2
\end{equation*}
Moreover, let
${\mcl{V}}_p\{{\mcl{Q}}_p, {\mcl{S}}_p, {\mcl{R}}_p\}$
be parameterized as,
\begin{equation*}
{\mcl{V}}_p =
\bmat{
{\mcl{Q}}_p & {\mcl{S}}_p \\
{\mcl{S}}_p^* & {\mcl{R}}_p
}
\end{equation*}
\end{definition}
In this framework, the multipliers $\mcl M_{\mcl V}$ and
$\mcl M_{\hat{\mcl{V}}}$ serve distinct but complementary roles.
The multiplier $\mcl M_{\mcl V}$ characterizes the behavior of uncertainties,
nonlinearities, and other perturbations captured by the operator $\Delta$,
ensuring that the IQC is satisfied for the class of admissible disturbances.
In contrast, the multiplier $\mcl M_{\hat{\mcl{V}}}$ encodes performance
requirements for the nominal system, such as $L_2$-gain bounds, passivity
conditions, or other desired input-output properties.
Together, these multipliers enable simultaneous verification of robust stability
against perturbations and achievement of specified performance objectives.
% The interconnected structure for robust stability and performance analysis
% can be viewed in Figure~\ref{fig:block structure}.

\subsection{Theorem for joint robust stability and performance}
We now present the main theoretical result that provides conditions for
simultaneous robust stability and performance verification. For brevity we will refer to the hard IQC as,
\begin{equation}
    \ip{p_T \Psi\bmat{z \\ w}}{p_T \mcl{M}_\mcl{V}\Psi\bmat{z \\ w}}_{\mbf L_{[a, b]}} \leq -\epsilon \|p_T w\|^2_{\mbf L_{[a, b]}}
\end{equation}
Let $\Psi_\Delta = \bmat{\Psi_{\Delta,11} & \Psi_{\Delta,12} \\ \Psi_{\Delta,21} & \Psi_{\Delta,22}}$ and define $\Psi_{ij} =
 \Psi_{\Delta,ij}\oplus \delta_{ij}I,\quad i,j\in\{1,2\}$, then $\Psi = \bmat{\Psi_{11} & \Psi_{12} \\ \Psi_{21} & \Psi_{22}}$.
Let $w=\bmat{w_\Delta \\ w_p}$ and $z=\bmat{z_\Delta \\ z_p} = Pw$.
Moreover, define $\mcl{Q} = \mcl{Q}_\Delta \oplus \mcl{Q}_p$,
$\mcl{S} = \mcl{S}_\Delta \oplus \mcl{S}_p$, and
$\mcl{R} = \mcl{R}_\Delta\oplus\mcl{R}_p$, then
$\mcl{V}\{\mcl{Q}, \mcl{S}, \mcl{R}\}$.

\begin{theorem}[Robust Stability and Performance]\label{thm:robust-stability-performance}
Let $P$, $\Delta\in\mbf{\Delta}$ be given, and there exists $\Psi_\Delta$,
$\mcl{M}_{\mcl{V}_\Delta}$, such that,
\begin{itemize}
\item The interconnection $P\star\Delta$ is well--posed for all
$\Delta\in\mbf{\Delta}$
\item The IQC in Definition~\ref{def:uncertainty-iqc} is satisfied for all
$\Delta\in\mbf{\Delta}$.
\end{itemize}
Then the interconnection $P \star \Delta$ is \textbf{robustly stable} and
\textbf{robust performance} is achieved with respect to some
$\mcl{M}_{{\mcl{V}}_p}$, in Definition~\ref{def:robust-performance}, with
$\mcl{Q}_p \succeq 0$, if,
\begin{equation}\label{eq:robust-stability-performance-condition}
\begin{gathered}
\ip{p_T\Psi\bmat{z\\w}}
{p_T\mcl{M}_{\mcl V}\Psi\bmat{z\\w}}_{\mbf L_{[a, b]}} \\
\leq
-\epsilon
\|p_Tw\|^2_{\mbf L_{[a, b]}}
\end{gathered}
\end{equation}
for all $w_\Delta \in \mbf L_{e,[a,b]}^{\mbf{n}_{w_\Delta}}$,
$w_p\in \mbf L_{e,[a,b]}^{\mbf{n}_{w}}$ and $T \geq 0$.
\end{theorem}
\begin{proof}
See Appendix~\ref{proof:robust-stability-performance}
\end{proof}
Theorem~\ref{thm:robust-stability-performance} establishes a framework for deriving Linear PI
Inequalities (LPIs), which are used in assessing the robust performance of a
system.
